We trained Gemini models using TPUv5e and TPUv4 \cite{tpuv4}, depending on their sizes and configuration. Training Gemini Ultra used a large fleet of TPUv4 accelerators owned by Google across multiple datacenters. This represents a significant increase in scale over our prior flagship model PaLM-2 which presented new infrastructure challenges. Scaling up the number of accelerators results in a proportionate decrease in the mean time between failure of hardware in the overall system. We minimized the rate of planned reschedules and preemptions, but genuine machine failures are commonplace across all hardware accelerators at such large scales.

TPUv4 accelerators are deployed in ``SuperPods'' of 4096 chips, each connected to a dedicated optical switch, which can dynamically reconfigure 4x4x4 chip cubes into arbitrary 3D torus topologies in around 10 seconds \citep{tpuv4}. For Gemini Ultra, we decided to retain a small number of cubes per superpod to allow for hot standbys and rolling maintenance. 

TPU accelerators primarily communicate over the high speed inter-chip-interconnect, but at Gemini Ultra scale, we combine SuperPods in multiple datacenters using Google's intra-cluster and inter-cluster network \cite{jupiter,b4,pony}. Google's network latencies and bandwidths are sufficient to support the commonly used synchronous training paradigm, exploiting model parallelism within superpods and data-parallelism across superpods.

The `single controller' programming model of Jax~\cite{jax} and Pathways~\cite{pathways} allows a single Python process to orchestrate the entire training run, dramatically simplifying the development workflow. The GSPMD partitioner~\cite{gspmd} in the XLA compiler partitions the training step computation, and  the MegaScale XLA compiler~\cite{xla} pass statically schedules appropriate collectives so that they maximally overlap with the computation with very little variation in step time.

Maintaining a high goodput\footnote{We define goodput as the time spent computing useful new steps over the elapsed time of the training job.} at this scale would have been impossible using the conventional approach of periodic checkpointing of weights to persistent cluster storage. For Gemini models, we instead made use of redundant in-memory copies of the model state, and on any unplanned hardware failures, we rapidly recover directly from an intact model replica. Compared to both PaLM and PaLM-2 \cite{palm2}, this provided a substantial speedup in recovery time, despite the significantly larger training resources being used. As a result, the overall goodput for the largest-scale training job increased from 85\% to 97\%.

Training at unprecedented scale invariably surfaces new and interesting systems failure modes - and in this instance one of the problems that we needed to address was that of ``Silent Data Corruption (SDC)''~\cite{sdcs, sdc2, hochschild2021cores}. Although these are extremely rare, the scale of Gemini models means that we can expect SDC events to impact training every week or two. Rapidly detecting and removing faulty hardware required several new techniques that exploit deterministic replay to isolate incorrect computations, combined with proactive SDC scanners on idle machines and hot standbys. Our fully deterministic infrastructure allowed us to quickly identify root causes (including hardware failures) during the development leading up to the Ultra model, and this was a crucial ingredient towards stable training.
